\subsection{Incidence evaluation of the terminal chart}
\label{subsec:k-incidencias-residentes}

The full state retains orientation, sheet, carry, boundary signature,
memory, and the dodecaphase ledger. On its compatible sublattice, the signed
reader is the projection of the generated coordinate and agrees with
\[
 R_{\mathrm{sgn}}=\operatorname{pr}_{U}
 =\Pi_H\circ\operatorname{pr}_{C}
 =\Pi_H\circ\mathcal E_{108}^{90,120}\circ\mathcal R_{12}.
\]
These identities act on the enriched state, not on the raw projection of
the one hundred and eight steps that forgets those fields. Descent through
that projection would require constancy on its fibres and is not assumed here.

To make the incidence presentation explicit, let \(\mathscr L\) be the
oriented ledger of the full state. Each visit retains its route, APP
position, sheet, time, multiplicity, and boundary incidence. Evaluation
on the sum or product sheet determines \(n_v=r_v+9q_v\), with
\(1\leq r_v\leq9\); the residue is not an independent input. For
\(E_m=\{9(m-1)+1,\ldots,9m\}\), define
\[
 \mathscr V_m=\{v:t(v)\in E_m\},\qquad
 \mathscr T_m=\{\theta:\operatorname{origen}(\theta)\in E_m\},
\]
and the counts
\[
\begin{aligned}
 A_m&=\sum_{v\in\mathscr V_m}w(v)\mathbf1_{\{1,3,5,7\}}(r(v)),\\
 C_m&=-\sum_{\theta\in\mathscr T_m}u(\theta)\sigma(\theta)
     =\sum_{j\geq0}\bigl(\nu^-_{120,m}(j)-\nu^+_{120,m}(j)\bigr),\\
 V_m&=\sum_{v\in\mathscr V_m}w(v)\mathbf1_{\{9\}}(r(v))
                         \epsilon_\partial(v).
\end{aligned}
\]
Here \(w\) and \(u\) are visit and trace multiplicities,
\(\sigma\) retains orientation, and \(\epsilon_\partial\) is \(1\) on
the corona and \(-1\) in the interior. The traces in \(\nu^\pm\) are
enumerated by the ledger rule; the integer sum requires finite support or
a construction that determines its meaning. The number of chronological
incidences and the reduced length \(120(1+3j)\) are distinct quantities:
their relation retains the reader's unit of length.

With \([x]_{10}\in\{0,\ldots,9\}\), the quotients in
\(x=[x]_{10}+10\lfloor x/10\rfloor\) are also retained. The decimal
readout and transverse extraction are
\[
\begin{aligned}
 z_m&=100[A_m]_{10}+10[C_m]_{10}+[V_m]_{10},\\
 \mathcal E_{\mathrm{cnt}}(\mathscr L)
   &=\bigl((z_m-z_{m+3})_{m=1}^{12},
           (z_m-z_{m+4})_{m=1}^{12},\sum_{m=1}^{12}z_m\bigr).
\end{aligned}
\]
Indices are cyclic modulo twelve. Thus \(\mathcal E_{\mathrm{cnt}}\)
is the presentation \((D_3z,D_4z,Q(z))\) of the extractor, with the same
convention as in \S\ref{subsec:k-transversal}. Its arguments are the
ledger counts, not \(K\), \(U\), alpha, or expected transverse coordinates.
The count presentation applies when the complete ledger is unfolded.
It does not assert that the ledger is recoverable after orientation and
memory have been forgotten. The harmonic evaluation and orbital inversion
developed below retain the compatible domain of the terminal chart.
